\subsection{Motivation}

Unfortunately and puzzlingly the first attempts to consider space-like~\cite{Craps:2002:StringPropagationPresence} or light-like singularities~\cite{Liu:2002:StringsTimeDependentOrbifold,Liu:2002:StringsTimeDependentOrbifolds} by means of orbifold  techniques yielded divergent four points \emph{closed string} amplitudes (see \cite{Cornalba:2004:TimedependentOrbifoldsString,Craps:2006:BigBangModels} for reviews).

These singularities are commonly assumed to be connected to a large backreaction of the incoming matter into the singularity due to the exchange of a single graviton~\cite{Berkooz:2003:CommentsCosmologicalSingularities,Horowitz:2002:InstabilitySpacelikeNull}.
This claim was already questioned in the literature where the $O$-plane orbifold was constructed.
This orbifold should in fact be stable against the gravitational collapse but it exhibits divergences in the amplitudes (see the discussion in \cite{Cornalba:2004:TimedependentOrbifoldsString}).
In what follows we show a direct computation showing that the presence of the divergence is not related to a gravitational response.

What has gone unnoticed is that in the Null Boost Orbifold (\nbo) \cite{Liu:2002:StringsTimeDependentOrbifold} even the four \emph{open string} tachyons amplitude is divergent.
Since we are working at tree level gravity is not an issue.
In fact in Equation (6.16) of \cite{Liu:2002:StringsTimeDependentOrbifold} the four tachyons amplitude in the divergent region reads
\begin{equation}
  A_4 \sim \int\limits_{q \sim \infty} \frac{\dd{q}}{\abs{q}} \ccA( q )
\end{equation}
where $\ccA_{\text{closed}}( q ) \sim q^{4 - \ap \norm{\vb{p}_{\perp}}^2}$ and $\ccA_{\text{closed}}( q ) \sim q^{1 - \ap \norm{\vb{p}_{\perp}}^2} \tr\qty( \liebraket{T_1}{T_2}_+ \liebraket{T_3}{T_4}_+ )$ ($T_i$ for $i = 1,\, 2,\, 3,\, 4$ are Chan-Paton matrices).
Moreover divergences in string amplitudes are not limited to four points: interestingly we show that the open string three point amplitude with two tachyons and the first massive state may be divergent when some \emph{physical} polarisations are chosen.
The true problem is therefore not related to a gravitational issue but to the non existence of the effective field theory.
In fact when we express the theory using the eigenmodes of the kinetic terms some coefficients do not exist, not even as a distribution.
This holds true for both open and closed string sectors since it manifests also in the four scalar contact term.
The issue can be roughly traced back to the vanishing volume of a subspace and the existence of a discrete zero mode of the Laplacian on this subspace.

As an introduction to the problem we first deal with singularities of the open string sector.
We try to build a consistent scalar \qed and show that the vertex with four scalar fields is ill defined.
Divergences in scalar QED are due to the behaviour of the eigenfunctions of the scalar d'Alembertian near the singularity but in a somehow unexpected way.
Near the singularity $u = 0$ in lightcone coordinates almost all eigenfunctions behave as $\frac{1}{\sqrt{\abs{u}}} e^{i \frac{\cA}{u}}$ with $\cA \neq 0$.
The product of $N$ eigenfunctions gives a singularity $\abs{u}^{-N/2}$ which is technically not integrable.
However the exponential term $e^{i \frac{\cA}{u}}$ allows for an interpretation as distribution when $\cA = 0$ is not an isolated point.
When $\cA = 0$ is isolated the singularity is definitely not integrable and there is no obvious interpretation as a distribution.
Specifically in the \nbo we find $\cA \sim \frac{l^2}{k_+}$ where $l$ is the momentum along the compact direction.
As a consequence we find the eigenfunction associated to the discrete momentum $l = 0$ along the orbifold compact direction with an isolated $\cA = 0$.
It is the eigenfunction which is constant along that direction and it is the root of all divergences.

We then check whether the most obvious ways of regularizing the theory by making $\cA$ not vanishing may work.
The first regularisation we try is to use a Wilson line along the compact direction even though the diverging three point string amplitude involves an anti-commutator of the Chan-Paton factor therefore it is divergent also for a neutral string, i.e.\ for a string with both ends attached to the same D-brane.
This kind of string does not feel Wilson lines.
Moreover anti-commutators are present in amplitudes with massive states in unoriented and supersymmetric strings and therefore neither worldsheet parity nor supersymmetry can help.
The second obvious regularisation is the introduction of higher derivatives couplings to the Ricci tensor which is the only non vanishing tensor associated to the (regularised) metric.
In any case it seems that a sensible regularisation must couple to all open string in the same way and this suggests a gravitational coupling.
We then give a cursory look to whether closed string winding modes could help~\cite{Berkooz:2003:StringsElectricField}, as already suggested in~\cite{Liu:2002:StringsTimeDependentOrbifolds,Craps:2002:StringPropagationPresence} in analogy to the resolution of static singularities.
Twisted closed strings become massless near the singularity and they should in some way be included.
They generate a background potential $B_{\mu\nu}$ which is equivalent to a electromagnetic background from the open string perspective. 
Under a plausible modification of the scalar action which is suggested by the two-tachyons---two-photons amplitude the problems seem to be solvable.

In any case the origin of the string divergence seems to originate from the lack of contact terms in the effective field theory.
Since these terms arise from string theory also through the exchange of massive string states we examine three point amplitudes with one massive state.
A deeper understanding of the subject requires the study of the polarisations of the massive state on the orbifold as seen from the covering Minkowski space before the computation of the overlap of the wave functions.
We then go back to string theory and we verify that in the \nbo the open string three points amplitude with two tachyons and one first level massive string state does indeed diverge when some physical polarisation are chosen.

We then introduce the Generalized Null Boost Orbifold (\gnbo) as a generalization of the \nbo which still has a light-like singularity and is generated by one Killing vector.
However in this model there are two directions associated with $\cA$, one compact and one non compact.
We can then construct the scalar \qed and the effective field theory which extends it with the inclusion of higher order terms since all terms have a distributional interpretation~\cite{Estrada:2012:GeneralIntegral}.
However if a second Killing vector is used to compactify the formerly non compact direction, the theory has again the same problems as in the \nbo.
In the literature there are however also other attempts at regularizing the \nbo such as the Null Brane.
This kind of orbifold was originally defined in \cite{Figueroa-OFarrill:2001:GeneralisedSupersymmetricFluxbranes,Cornalba:2004:TimedependentOrbifoldsString} and studied in perturbation theory in \cite{Liu:2002:StringsTimeDependentOrbifolds}.
The Null Brane shares with the \gnbo the existence of a non compact direction on the orbifold.
In this case it is indeed possible to build single particle wave functions which leads to the convergence of the smeared amplitudes.

We finally present also a brief examination of the Boost Orbifold (\bo) where the divergences are generally milder~\cite{Horowitz:1991:SingularStringSolutions,Khoury:2002:BigCrunchBig}.
The scalar eigenfunctions behave in time $t$ as $\abs{t}^{\pm i\, \frac{l}{\Delta}}$ near the singularity but there is one eigenfunction which behaves as $\log(\abs{t})$ and again it is the constant eigenfunction along the compact direction which is the origin of all divergences.
In particular the scalar \qed on the \bo can be defined and the first term which gives a divergent contribution is of the form $\abs{\phi~\dphi}^2$, i.e.\ divergences are hidden into the derivative expansion of the effective field theory.
Again three points open string amplitudes with one massive state diverge.

\subsection{Scalar QED on NBO and Divergences} 
\label{sect:NOscalarQED}

As discussed the four open string tachyons amplitude diverges in the \nbo.
The literature on the subject (see for instance~\cite{Cornalba:2004:TimedependentOrbifoldsString} and references therein) suggests that this can be cured by the eikonal resummation.
We therefore consider the scalar \qed on the \nbo as a first approach.
In this case all eigenmodes can be written using elementary functions thus making the issues even more evident.
Its action is given by
\begin{equation}
  \rS_{\text{s}\qed}
  =
  \int\limits_{\Omega} \dd[D]{x}\,
  \sqrt{- \det g}
  \qty(
    - \qty(D^{\mu} \phi)^*\, D_{\mu} \phi
    - M^2  \qty(\phi^*)\, \phi
    - \frac{1}{4} f^{\mu\nu}\, f_{\mu\nu}
    - \frac{g_4}{4} \abs{\phi}^4
  ),
\end{equation}
with
\begin{equation}
  D_{\mu} \phi
  =
  \qty(\ipd{\mu} -i\, e\,  a_{\mu}) \phi,
  \qquad
  f_{\mu\nu}
  =
  \ipd{\mu} a_{\nu} - \ipd{\nu} a_{\mu}.
\end{equation}
We reserve small letters for quantities defined on the orbifold and capital letters for those defined in flat space.
Moreover $\Omega$ denotes the orbifold.
We will construct directly both the scalar and the spin-1 eigenfunctions  which we can use as a starting point for the perturbative computations.


\subsubsection{Geometric Preliminaries}
\label{sec:geometric_preliminaries_nbo}

In Minkowski spacetime $\ccM^{1,D-1}$ with coordinates $\qty(x^{\mu}) = \qty(x^+,\, x^-,\, x^2,\, \vb{x})$
and metric
\begin{equation}
  \dss[2]{s}
  =
  - 2 \dd{x^+} \dd{x^-}
  + \qty(\dd{x^2})^2
  + \eta_{ij} \dd{x}^i \dd{x}^j,
\end{equation}
we consider the following change of coordinates to $\qty(x^{\alpha}) = (u,\, v,\, z,\, \vb{x})$
\begin{equation}
  \begin{cases}
    x^- & = u
    \\
    x^2 & = \Delta u z
    \\
    x^+ & = v + \frac{1}{2} \Delta^2 u z^2
  \end{cases}
  \qquad
  \Leftrightarrow
  \qquad
  \begin{cases}
    u & = x^-
    \\
    z & = \frac{x^2}{\Delta\, x^-}
    \\
    v & = x^+ - \frac{1}{2} \frac{(x^2)^2}{x^-}
  \end{cases}.
  \label{eq:NBO_coordinates}
\end{equation}
Then the metric becomes:
\begin{equation}
  \dss[2]{s}
  =
  - 2\, \dd{u}\, \dd{v}
  + \qty(\Delta u )^2 (\dd{z})^2
  + \eta_{ij} \dd{x}^i \dd{x}^j,
\end{equation}
along with the non vanishing geometrical quantities
\begin{equation}
  -\det g = \qty( \Delta u )^2,
\end{equation}
and
\begin{equation}
  \tensor{\Gamma}{_z^v_z} = \Delta^2 u,
  \qquad
  \tensor{\Gamma}{_u^z_z} = u^{-1}.
\end{equation}
Riemann and Ricci tensor components however vanish since at this stage we only performed a change of coordinates from the original Minkowski spacetime.
Locally it is the same as the \nbo and they must have the same local differential geometry.

The \nbo is introduced by identifying points along the orbits of the Killing vector:
\begin{equation}
  \begin{split}
    \kappa
    & =
    - i \qty(2 \pi \Delta) J_{+2}
    \\
    & =
    \qty(2 \pi \Delta)\, (x^2 \ipd{+} + x^- \ipd{2})
    \\
    & =
    2 \pi \ipd{z},
  \end{split}
  \label{eq:nbo_killing_vector}
\end{equation}
in such a way that
\begin{equation}
  x^{\mu} \equiv \cK^{n}\, x^{\mu},
  \qquad
  n \in \Z,
\end{equation}
where $\cK^{n}= e^{n\kappa}$, leads to the identifications
\begin{equation}
  x=
  \mqty( x^- \\  x^2 \\ x^+ \\ \vb{x} )
  \equiv
  \cK^{n} x
  =
  \mqty(%
    x^- \\
    x^2 + n \qty(2 \pi \Delta) x^- \\
    x^+ + n \qty(2 \pi \Delta) x^2 + \frac{1}{2} n^2 \qty(2 \pi \Delta)^2 x^- \\
   \vb{x}
  )
\end{equation}
or to 
\begin{equation}
  \qty( u,\, v,\, z,\, \vb{x} )
  \equiv
  \qty( u,\, v,\, z + 2 \pi n,\, \vb{x} )
\end{equation}
in coordinates $\qty(x^{\alpha})$ where $\kappa = 2 \pi \ipd{z}$ is a global Killing vector.

As a reference for the future, we notice that we could regularise the metric as
\begin{equation}
  \dss[2]{s}
  =
  - 2\, \dd{u}\, \dd{v}
  + \Delta^2 \qty(u^2 + \epsilon^2) (\dd{z})^2
  + \eta_{ij} \dd{x}^i \dd{x}^j.
\end{equation}
The non vanishing geometrical quantities are then:
\begin{equation}
  -\det g = \Delta^2 \qty(u^2 + \epsilon^2),
\end{equation}
and
\begin{equation}
  \tensor{\Gamma}{_z^v_z} = \Delta^2 u,
  \qquad
  \tensor{\Gamma}{_u^z_z} = \frac{u}{u^2 + \epsilon^2},
\end{equation}
which lead to the following Riemann and Ricci tensor components:
\begin{equation}
  \tensor{R}{^z_u_z_u} = - \frac{\epsilon^2}{\qty(u^2+ \epsilon^2)^2},
  \quad
  \tensor{R}{^v_z_z_u} = - \frac{\Delta^2 \epsilon^2}{u^2 + \epsilon^2},
  \quad
  \tensor{Ric}{_u_u} = -\frac{\epsilon^2}{\qty(u^2+ \epsilon^2)^2}.
\end{equation}
Since $\delta_{\text{reg}}(u) = \frac{1}{\pi} \frac{\epsilon}{u^2+ \epsilon^2}$ then $\tensor{R}{^z_u_z_u} = - \pi^2 \qty[ \delta_{\text{reg}}(u) ]^2$.


\subsubsection{Free Scalar Action} 
 
We study the eigenmodes of the Laplacian operator to diagonalize the scalar kinetic term given by:\footnotemark{}
\footnotetext{%
  The factor $-g^{\alpha\beta}$ is due to the choice of the East coast convention for the metric, namely:
  \begin{equation*}
    - g^{\alpha\beta}
    \ipd{\alpha} \phi^*\, \ipd{\beta} \phi
    -
    M^2 \phi^*\, \phi
    \sim
    \abs{\dot{\phi}}^2 - M^2 \abs{\phi}^2
    \sim
    \rE^2 - M^2.
    \end{equation*}
}
\begin{equation}
  \begin{split}
    \rS_{\text{s}\qed}^{(\text{kinetic})}\qty[ \phi ]
    & =
    \int\limits_{\Omega} \dd[D]{x}\,
    \sqrt{- \det g}~
    \qty(%
      - g^{\alpha\beta} \ipd{\alpha} \phi^*\, \ipd{\beta} \phi
      - M^2  \phi^*\, \phi
    )
    \\
    & =
    \int \dd[D-3]{\vb{x}}\,
    \int \dd{u}\,
    \int \dd{v}\,
    \finiteint{z}{0}{2\pi}
    \abs{\Delta u}
    \\
    & \times
    \qty(%
      \ipd{u} \phi^*\, \ipd{v} \phi\,
      +
      \ipd{v} \phi^*\, \ipd{u} \phi\,
      -
      \frac{1}{\qty(\Delta u)^2} \ipd{z} \phi^*\, \ipd{z} \phi\,
      -
      \ipd{i} \phi^*\, \ipd{i} \phi
      -
      M^2 \phi^*\, \phi
    ).
  \end{split}
\end{equation}
The solution to the equation of motion is enough when we want to perform the canonical quantization.
Since we use Feynman diagrams we consider the path integral approach: we take off-shell modes and solve the eigenvalue problem $\square \phi_r = r \phi_r$.
Comparing with the flat case we see that $r$ is $2\, k_-\, k_+ - \norm{\vb{k}}^2$ when $k$ is the impulse in flat coordinates.
We therefore have
\begin{equation}
  -2 \ipd{u} \ipd{v} \phi_r
  -
  \frac{1}{u} \ipd{v} \phi_r
  +
  \frac{1}{\qty(\Delta u)^2} \ipd{z}^2 \phi_r
  +
  \ipd{i}^2 \phi_r
  =
  r \phi_r.
  \label{eq:nbo_eom}
\end{equation}
Using Fourier transforms it follows that the eigenmodes are
\begin{equation}
  \phi_{\kmkr}\qty(u,\, v,\, z,\, \vb{x})
  =
  e^{i k_+ v + i l z + i \vb{k} \cdot \vb{x}}\,
  \tphi_{\kmkr}(u),
\end{equation}
with
\begin{equation}
  \tphi_{\kmkr}(u)
  =
  \frac{1}{\sqrt{\qty( 2 \pi )^D~ \abs{2 \Delta k_+\, u}}}
  e^{
    - i \frac{l^2}{2 \Delta^2 k_+} \frac{1}{u}
    + i \frac{\norm{\vb{k}}^2 + r}{2 k_+} u
  },
\end{equation}
and
\begin{equation}
  \phi^*_{\kmkr}\qty(u,\, v,\, z,\, \vb{x})
  =
  \phi_{\mkmkr}\qty(u,\, v,\, z,\, \vb{x}).
\end{equation}
We chose the numeric factor in order to get a canonical normalisation:
\begin{equation}
  \begin{split}
    &
    \qty( \phi_{\kmkrN{1}},\,  \phi_{\kmkrN{2}} )
    \\
    = &
    \int \dd[D-1]{\vb{x}}\,
    \int \dd{u}\, \int \dd{v}\,
    \finiteint{z}{0}{2\pi}
    \abs{\Delta u}\,
    \phi_{\kmkrN{1}}\, \phi_{\kmkrN{2}}
    \\
    = &
    \delta^{D-3}( \vb{k}_{\qty(1)} + \vb{k}_{\qty(2)})\,
    \delta( r_{(1)} - r_{(2)})\,
    \delta( k_{\qty(1)\, +} + k_{\qty(2)\, +})\,
    \delta_{l_{\qty(1)} + l_{\qty(2)},\, 0}.
  \end{split}
\end{equation} 
We can then perform the off-shell expansion
\begin{equation}
  \phi\qty(u,\, v,\, z,\, \vb{x})
  =
  \int \dd[D-3]{\vb{k}}
  \int \dd{k_+}
  \int \dd{r}
  \infinfsum{l}
  \cA_{\kmkr}\,
  \phi_{\kmkr}\qty(u,\, v,\, z,\, \vb{x}),
\end{equation}
such that the scalar kinetic term becomes 
\begin{equation}
  \rS_{\text{s}\qed}^{(\text{kinetic})}\qty[ \cA ]
  =
  \int \dd[D-3]{\vb{k}}\,
  \int \dd{k_+} 
  \int \dd{r}
  \infinfsum{l}
  \qty(r - M^2)\,
  \cA_{\kmkr}\,
  \cA_{\kmkr}^*.
\end{equation}


\subsubsection{Free Photon Action}

The action of the free photon can be written as
\begin{align}
  \rS_{\text{s}\qed}^{(\text{kinetic})}\qty[ a ]
  =
  \int\limits_{\Omega} \dd[D]{x}\,
  \sqrt{-\det g}\,
  \qty(%
   - \frac{1}{2} g^{\alpha\beta} g^{\gamma\delta}
   D_{\alpha} a_{\gamma} \qty( D_{\beta} a_{\delta} - D_{\delta} a_{\beta})
  ).
\end{align}
We choose to enforce the Lorenz gauge:\footnotemark{}
\footnotetext{%
  Indeed it is exactly the usual Lorenz gauge since locally the space is Minkowski.
}
\begin{equation}
  D^{\alpha} a_{\alpha}
  =
  - \frac{1}{u} a_{v}
  - \ipd{u} a_{v}
  - \ipd{v} a_{u}
  + \frac{1}{\Delta^2 u^2} \ipd{z} a_z
  + \eta^{ij} \ipd{i} a_j
  =
  0.
  \label{eq:Lorenz_gauge}
\end{equation}
As covariant derivatives commute since we are locally flat, the \eom read $\qty(\square a)_{\alpha} = 0$.
Explicitly we have:
\begin{equation}
  \begin{split}
    \qty( \square a )_u
    & =
    \frac{1}{u^2} a_{v}
    -
    \frac{2}{\Delta^2 u^3} \ipd{z} a_z
    +
    \qty[
      -
      2 \ipd{u} \ipd{v}
      -
      \frac{1}{u} \ipd{v}
      +
      \frac{1}{\Delta^2 u^2} \ipd{z}^2
      +
      \eta^{ij} \ipd{i} \partial_j
    ]
    a_u,
    \\
    \qty( \square a )_v
    & =
    \qty[
    -
      2 \ipd{u} \ipd{v}
      -
      \frac{1}{u} \ipd{v}
      +
      \frac{1}{\Delta^2 u^2} \ipd{z}^2
      +
      \eta^{ij} \ipd{i} \partial_j
    ]
    a_v,
    \\
    \qty( \square a )_z
    & =
    -
    \frac{2}{u} \ipd{z} a_v
    +
    \qty[
      -
      2 \ipd{u} \ipd{v}
      +
      \frac{1}{u} \ipd{v}
      +
      \frac{1}{\Delta^2 u^2} \ipd{z}^2
      +
      \eta^{ij} \ipd{i} \partial_j
    ]
    a_z,
    \\
    \qty( \square a )_i
    & =
    \qty[
      -
      2 \ipd{u} \ipd{v}
      -
      \frac{1}{u} \ipd{v}
      +
      \frac{1}{\Delta^2 u^2} \ipd{z}^2
      +
      \eta^{ij} \ipd{i} \partial_j
    ]
    a_i.
  \end{split}
\end{equation}

As in the previous scalar case we are actually interested in solving
the eigenmodes problem $\qty(\square a)_\alpha= r \,a_\alpha$.
We proceed hierarchically: first we solve for $a_v$ and $a_i$ whose equations are the same as in the scalar field, then we insert the solutions as a source in the equation for $a_z$ and eventually we solve for $a_u$.\footnotemark{}
\footnotetext{%
  Notice that inside the square brackets of the differential equation for $a_z$ there is a different sign for the term $\frac{1}{u} \ipd{v}$ with respect to the equation for the scalar field.
}
We get the solutions:
\begin{equation}
  \begin{split}
    \norm{\tildea_{\kmkr\, \alpha}(u)}
    \,=
    \mqty(%
      \tildea_u
      \\
      \tildea_v
      \\
      \tildea_z
      \\
      \tildea_i
    )
    & =
    \sum\limits_{%
      \underline{\alpha}
      \in 
      \qty{ \underline{u}, \underline{v}, \underline{z},\underline{i} }
    }
    \cE_{\kmkr\, \underline{\alpha}}
    \norm{\tildea^{\underline{\alpha}}_{\kmkr\, \alpha}(u)}
    \\
    & =
    \cE_{\kmkr\, \underline{u}}
    \mqty(
      1
      \\
      0
      \\
      0
      \\
      0
    )\,
    \tphi_{\kmkr}(u)
    \\
    & +
    \cE_{\kmkr\, \underline{v}}
    \mqty(
      \frac{i}{2 k_+ u}
      +
      \frac{1}{2} \qty( \frac{l}{\Delta k_+} )^2 \frac{1}{u^2}
      \\
      1
      \\
      \frac{l}{k_+}
      \\
      0
    )\,
    \tphi_{\kmkr}(u)
    \\
    & +
    \cE_{\kmkr\, \underline{z}}
    \mqty(
      \frac{l}{\Delta k_+ \abs{u}}
      \\
      0
      \\
      \Delta \abs{u}
      \\
      0
    )\,
    \tphi_{\kmkr}(u)
    \\
    & +
    \cE_{\kmkr\, \underline{j}}
    \mqty(
      0
      \\
      0
      \\
      0
      \\
      \delta_{\underline{ij}}
    )\,
    \tphi_{\kmkr}(u),
    \label{eq:Orbifold_spin1_pol}
  \end{split}
\end{equation}
then we can expand the off-shell fields as
\begin{equation}
  a_{\alpha}\qty(u,\, v,\, z,\, \vb{x} )
  =
  \int \ccD k
  \sum\limits_{%
    \underline{\alpha}
    \in
    \qty{ \underline{u}, \underline{v}, \underline{z},\underline{i} }
  }
  \infinfsum{l}
  \cE_{\kmkr\, \underline{\alpha}}\,
  {a}^{\underline{\alpha}}_{\kmkr\, \alpha}\qty(u,\, v,\, z,\, \vb{x} ),
\end{equation}
where ${a}^{\underline{\alpha}}_{\kmkr\, \alpha}\qty(u,\, v,\, z,\, \vb{x}) = \tildea^{\underline{\alpha}}_{\kmkr\, \alpha}(u)\, e^{i\, \qty( k_+ v + l z + \vb{k} \cdot \vb{x})}$ and $\int \ccD k = \int \dd[D-3]{\vb{k}} \int \dd{k_+} \int \dd{r}$.

We can also compute the normalisation as
\begin{equation}
  \begin{split}
    \qty(a_{(1)},\, a_{(2)})
    & =
    \int \dd[D-3]{\vb{x}}
    \int \dd{u}
    \int \dd{v}
    \finiteint{z}{0}{2\pi}
    \abs{\Delta u}
    \\
    & \times
    g^{\alpha\beta}\,
    a_{\kmkrN{1}\, \alpha}\,  a_{\kmkrN{2}\, \beta}
    \\
    & =
    \cE_{\kmkrN{1}} \circ \cE_{\kmkrN{2}}
    \\
    & \times
    \delta^{D-3}( \vb{k}_{\qty(1)} + \vb{k}_{\qty(2)})\,
    \delta( r_{(1)} - r_{(2)})\,
    \delta( k_{\qty(1)\, +} + k_{\qty(2)\, +})\,
    \delta_{l_{\qty(1)} + l_{\qty(2)},\, 0},
  \end{split}
\end{equation}
where:\footnotemark{}
\footnotetext{%
  We use a shortened version of the polarizations $\cE$ for the sake of readability.
  We write $\cE_{(n)\, \underline{\alpha}} = \cE_{\kmkrN{n}\, \underline{\alpha}}$ thus hiding the understood dependence of the components of $\cE_{(n)}$ on the momenta.
}
\begin{equation}
  \begin{split}
    \cE_{(1)} \circ \cE_{(2)}
    =
    - \cE_{(1)\, \underline{u}}\, \cE_{(2)\, \underline{v}}
    - \cE_{(1)\, \underline{v}}\, \cE_{(2)\, \underline{u}}
    + \cE_{(1)\, \underline{z}}\, \cE_{(2)\, \underline{z}}
    +
    \eta^{\underline{ij}}\,
    \cE_{(1)\, \underline{i}}\, \cE_{(2)\, \underline{j}}.
  \end{split}
\end{equation}
Finally the Lorenz gauge reads
\begin{equation}
   \eta^{i \underline{j}}\, k_i\, \cE_{\kmkr\, \underline{j}}
   -
   k_+\, \cE_{\kmkr\, \underline{u}}
   -
   \frac{\vb{k}^2 + r}{2\, k_+} \cE_{\kmkr\, \underline{v}}
   =
   0,
  \label{eq:explicit_orbifold_Lorenz}
\end{equation}
which does not impose any constraint on the transverse polarization
$\cE_{\kmkr\, \underline{z}}$.
The photon kinetic term becomes 
\begin{equation}
  \rS_{\text{s}\qed}^{(\text{kinetic})}\qty[ \cE ]
  =
  \int \dd[D-3]{\vb{k}}
  \int \dd{k_+}
  \int \dd{r}
  \infinfsum{l}\,
  \frac{r}{2}\,
  \cE_{\kmkr}\,
  \circ
  \cE_{\kmkr}^*.
\end{equation}


\subsubsection{Cubic Interaction}

With the definition of the d'Alembertian eigenmodes we can now examine the cubic vertex which reads
\begin{equation}
  \rS_{\text{s}\qed}^{(\text{cubic})}\qty[\phi,\, a]
  =
  \int\limits_{\Omega} \dd[D]{x}\,
  \sqrt{- \det g}\,
  \qty(%
    -i\, e\,
    g^{\alpha\beta}
    a_{\alpha}
    \qty(%
      \phi^*\, \ipd{\beta} \phi
      -
      \ipd{\beta} \phi^*\, \phi
    )
  ).
\end{equation}
Its computation involves integrals such as
\begin{equation}
  \int \dd{u}\,
  \abs{\Delta u}\,
  \qty(\frac{l}{u})^2
  \finiteprod{i}{1}{3} \tphi_{\kmkrN{i}}
  \sim
  \int_{u \sim 0} \dd{u}\, 
  \qty(\frac{l^2}{\abs{u}^{\frac{5}{2}}})
  e^{%
    -i \finitesum{i}{1}{3} \frac{l_{\qty(i)}^2}{2\, \Delta^2 k_{\qty(i)\, +)}}
    \frac{1}{u}
  },
\end{equation}
and
\begin{equation}
  \int \dd{u}\,
  \abs{\Delta u}\,
  \qty(\frac{1}{u})
  \finiteprod{i}{1}{3} \tphi_{\kmkrN{i}}
  \sim
  \int_{u \sim 0} \dd{u}\, 
  \qty(\frac{1}{u\, \abs{u}^{\frac{1}{2}}})
  e^{%
    -i \finitesum{i}{1}{3} \frac{l_{\qty(i)}^2}{2\, \Delta^2 k_{\qty(i)\, +}}
    \frac{1}{u}
  },
  \label{eq:nbo_div_integral}
\end{equation}
which can be interpreted as hints that the theory is troublesome.
The first integral diverges if the exponential functions are all equal to unity.
Fortunately it happens when all factors $l_{\qty(i)}$ (where $i = 1,\, 2,\, 3$) vanish.
In this case however the integral vanishes if we set $l_{\qty(i)} = 0$ before its evaluation.
This however suggests that when all $l_{\qty(i)} = 0$, i.e.\ when the eigenfunctions are constant along the compact direction $z$, something suspicious is happening.
On the other side when at least one $l$ is different from zero we have an integral such as:
\begin{equation}
  \int_{u \sim 0} \dd{u}\, 
  \abs{u}^{-\nu}\, e^{i \frac{\cA}{u}}
  \sim
  \int_{t \sim \infty} \dd{t}\, 
  t^{\nu-2}\, e^{i \cA t}.
\end{equation}
All $l_{\qty(i)}$ are discrete but $k_{\qty(i)\, +}$ are not thus $\cA$ has an isolated zero.
Otherwise it has continuous value and may be given a distributional meaning, similar to a derivative of the Dirac delta function.
The second integral has the same issues when all $l_{\qty(*)} = 0$ but, since it is not proportional to any $l$ as it stands, it is divergent unless we consider a principal part regularization.

We can give in any case meaning to the cubic terms
and we get:\footnotemark{}
\footnotetext{%
  The notation $(2) \rightarrow (3)$ meaning is that all previous terms inside the curly brackets appear again in exactly the same structure but with momenta of particle $(3)$ in place of those of particle $(2)$.
}
\begin{equation}
  \begin{split}
    \rS_{\text{s}\qed}^{(\text{cubic})}\qty[ \cA,\, \cE ]
    & =
    \finiteprod{i}{1}{3}
    \qty[%
      \int \dd[D-3]{\vb{k}_{\qty(i)}}\,
      \dd{r_{\qty(i)}}\,
      \dd{k_{\qty(i)\, +}}
      \sum_{l_{\qty(i)}}
    ]\,
    \qty(2\pi)^{D-1}
    \delta\qty(\finitesum{i}{1}{3} \vb{k}_{\qty(i)})\,
    \delta\qty(\finitesum{i}{1}{3} k_{\qty(i)\, +})\,
    \\
    & \times
    e~
    \delta_{l_{\qty(1)} + l_{\qty(2)} + l_{\qty(3)},\, 0}\,
    \qty(\cA_{\mkmkrN{2}})^*\, \cA_{\kmkrN{3}}
    \\
    & \times
    \left\lbrace
      \cE_{\kmkrN{1}\, \underline{u}}\,
      k_{\qty(2)\, +}\,
      \cI_{\qty{3}}^{\qty[0]}
    \right.
    \\
    & +
    \cE_{\kmkrN{1}\, \underline{z}}\,
    \frac{%
      k_{\qty(2)\, +} l_{\qty(1)}
      -
      l_{\qty(2)} k_{\qty(1)\, +}
    }{\Delta k_{\qty(1)\, +}}\,
    \cJ_{\qty{3}}^{\qty[-1]}
    \\
    & +
    \cE_{\kmkrN{1}\, \underline{v}}\,
    \ccF\qty(k_{\qty(1)\, +},\, l_{\qty(1)},\, k_{\qty(2)\, +},\, l_{\qty(2)},\, r_{\qty(2)},\, \vb{k}_{\qty(2)})
    \\
    & -
    \left.
      \eta^{\underline{i}\, j}\,
      \cE_{\kmkrN{1}\, \underline{i}}\,
      k_{(2)_j}\,
      \cI_{\qty{3}}^{\qty[0]}\,
      -
      \qty( (2) \rightarrow (3) )
    \right\rbrace,
  \label{eq:sQED_cubic_final}   
  \end{split}
\end{equation}
where
\begin{equation}
  \begin{split}
    \ccF\qty(k_{\qty(1)\, +},\, l_{\qty(1)},\, k_{\qty(2)\, +},\, l_{\qty(2)},\, r_{\qty(2)},\, \vb{k}_{\qty(2)})
    & =
    \frac{\norm{\vb{k}_{\qty(2)}}^2 + r_{\qty(2)}}{2\, k_{\qty(2)\, +}} \cI_{\qty{3}}^{\qty[0]}
    +
    i
    \frac{k_{\qty(2)\, +}}{2\, k_{\qty(1)\, +}}
    \cI_{\qty{3}}^{\qty[-1]}
    \\
    & +
    \frac{1}{2} \frac{k_{\qty(2)\, +}}{\Delta^2}
    \qty(%
      \frac{l_{\qty(1)}}{k_{\qty(1)\, +}}
      -
      \frac{l_{\qty(2)}}{k_{\qty(2)\, +}}
    )^2
    \cI_{\qty{3}}^{\qty[-2]}.
  \end{split}
\end{equation}
In the previous expressions we also defined for future use:
\begin{eqnarray}
  \cI_{(1) \dots (N)}^{\qty[\nu]}
  =
  \cI_{\qty{N}}^{\qty[\nu]}
  & = &
  \infinfint{u}\,
  \abs{\Delta u}\, u^{\nu}\,
  \finiteprod{i}{1}{N}
  \tphi_{\kmkrN{i}}
  \\
  \cJ_{\qty{N}}^{\qty[\nu]}
  & = &
  \infinfint{u}\,
  \abs{\Delta}\, \abs{u}^{1 + \nu}
  \finiteprod{i}{1}{N} \tphi_{\kmkrN{i}}.
\end{eqnarray}
For the sake of brevity from now on we use
\begin{eqnarray}
  \tphi_{\qty(i)} & = & \tphi_{\kmkrN{i}},
  \\
  \tphi_{\qty(i)} & = & \tphi_{\kmkrN{i}}
\end{eqnarray}
when not causing confusion.


\subsubsection{Quartic Interactions and Divergences}

The issue with the divergent vertex is even more visible when considering the quartic terms:
\begin{equation}
  \rS_{\text{s}\qed}^{(\text{quartic})}\qty[ \phi,\, a ]
  =
  \int\limits_{\Omega} \dd[D]{x}\,
  \sqrt{- \det g}\,
  \qty(%
    e^2\, g^{\mu\nu}\, a_{\mu} a_{\nu}\, \abs{\phi}^2
    -
    \frac{g_4}{4}\abs{\phi}^4
  ),
\end{equation}
which can be expressed using the modes as:
\begin{equation}
  \begin{split}
    \rS_{\text{s}\qed}^{(\text{quartic})}\qty[ \phi,\, a ]
    & =
    \finiteprod{i}{1}{4}
    \qty[%
      \int \dd[D-3]{\vb{k}_{\qty(i)}}
      \dd{k_{\qty(i)\, +}}
      \dd{r_{\qty(i)}}
      \sum_{l_{\qty(i)}}
    ]\,
    \qty(2\pi)^{D-1}
    \\
    & \times
    \delta\qty( \finitesum{i}{1}{4} \vb{k}_{\qty(i)} )\,
    \delta\qty( \finitesum{i}{1}{4} k_{\qty(i)\, +} )\,
    \delta_{l_{\qty(1)} + l_{\qty(2)} + l_{\qty(3)} + l_{\qty(4)},\, 0}
    \\
    & \times
    \left\lbrace
      e^2\,
      \qty(\cA_{\mkmkrN{3}})^* \cA_{\kmkrN{4}}
    \right\rbrace
    \\
    & \times
    \left[
     \qty(\cE_{\kmkrN{1}} \circ \cE_{\kmkrN{2}})\,
     \cI_{\qty{4}}^{\qty[0]}
    \right.
    \\
    & -
    \frac{i}{2}
    \cE_{\kmkrN{1}\, \underline{v}}\, \cE_{\kmkrN{2}\, \underline{v}}
    \qty(%
      \frac{1 }{k_{\qty(2)\, +}}
      +
      \frac{1}{k_{\qty(1)\, +}}
    )\,
    \cI_{\qty{4}}^{\qty[-1]}
    \\
    & +
    \left.
      \frac{1}{2}\,
      \frac{\cE_{\kmkrN{1}\, \underline{v}} \cE_{\kmkrN{2}\, \underline{v}} }{\Delta^2}
      \qty(%
        \frac{l_{\qty(1)}}{k_{\qty(1)\, +}}
        -
        \frac{l_{\qty(2)}}{k_{\qty(2)\, +}}
      )^2\,
      \cI_{\qty{4}}^{\qty[-2]}
    \right]
    \\
    & -
    \left.
      \frac{g_4}{4}\,
      \ccA\qty(\qty{k_+,\, l,\, \vb{k},\, r})\,
      \cI_{\qty{4}}^{\qty[0]}
    \right\rbrace,
  \end{split}
\end{equation}
where
\begin{equation}
  \begin{split}
    \ccA\qty(\qty{k_+,\, l,\, \vb{k},\, r})
    & =
    \qty(\cA_{\mkmkrN{1}})^*\,
    \qty(\cA_{\mkmkrN{2}})^*\,
    \\
    & \times
    \cA_{\kmkrN{3}}\,
    \cA_{\kmkrN{4}}.
  \end{split}
\end{equation}
When setting $l_{\qty(*)} = 0$ all the surviving terms are divergent.
The explicit behaviour is $\cI_{\qty{4}}^{\qty[0]} \sim \int \dd{u}\, \abs{u}^{1 -4 \times \frac{1}{2}}$ and $\cI_{\qty{4}}^{\qty[-1]} \sim \int \dd{u}\, u^{-1}\, \abs{u}^{1 - 4 \times \frac{1}{2}}$ since $\eval{\tphi}_{l = 0} \sim \abs{u}^{-\frac{1}{2}}$.
Higher order terms in the effective field theory have even worse behaviour.
This makes the theory ill defined and the string theory which should give this effective theory ill defined too.


\subsubsection{Failure of Obvious Divergence Regularizations}
\label{sec:saving}

From the discussion in the previous section the origin of the divergences is the sector $l = 0$.
When $l = 0$ the highest order singularity of the Fourier transformed d'Alembertian equation vanishes.
Explicitly we have:
\begin{equation}
  A\, \ipd{u} \tphi_{\kmkr}
  +
  B(u)\, \tphi_{\kmkr}
  =
  A\, e^{-\int^u \frac{B(u)}{A} du}\, 
  \ipd{u} \qty[ e^{+\int^u \frac{B(u)}{A} du} \tphi_{\kmkr} ]
  =
  0,
\end{equation}
with
\begin{equation}
  A = \qty(-2\, i\, k_+),
  \qquad
  B(u)
  =
  -\qty(\vb{k}^2 + r)
  -
  i\, k_+\, \frac{1}{u}
  -
  \frac{l^2}{\Delta^2}\, \frac{1}{u^2}.
\end{equation}
This implies the absence of the oscillating factor $e^{i \frac{\cA}{u} }$ when $l$ vanishes.
It follows that any deformation which prevents the coefficient of the highest order singularity from vanishing will do the trick.

The first and easiest possibility is to add a Wilson line along $z$, i.e.\ $a = \theta \dd{z}$.
This shifts $l \rightarrow l - e\, \theta$ and regularises the scalar \qed.
Unfortunately this does not work in the string theory where Wilson lines on D25-branes are not felt by the neutral strings starting and ending on the same D-brane.
In fact not all interactions involve commutators of the Chan-Paton factors which vanish for neutral strings.
For instance the interaction of two tachyons with the first massive state involves an anti-commutator as we discuss later. 
The anti-commutators are present also in amplitudes of supersymmetric strings with massive states and therefore the issue is not solved by supersymmetry.

A second possibility is to include higher derivative couplings to curvature as natural in the string theory.
If we regularise the metric in a minimal way as shown at the end of~\Cref{sec:geometric_preliminaries_nbo}, only $\tensor{Ric}{_u_u}$ does not vanish.
We can introduce:
\begin{equation}
  \begin{split}
    &
    S_{\mathrm{HE}}^{(\text{higher R})}\qty[ \phi,\, g ]
    \\
    = &
    \int\limits_{\Omega} \dd[D]{x}\,
    \sqrt{- \det g}\,
    \qty(%
      \finitesum{k}{1}{+\infty}
      \qty(\ap)^{2 k-1}\,
      \finiteprod{j}{1}{k}\,
      g^{\mu_j \nu_j}\, g^{\rho_j\sigma_j}\,
      \tensor{Ric}{_{\mu_j}_{\rho_j}}\,
      \qty(%
        \finitesum{s}{0}{2k}
        c_{k s}\, \ipd{\nu_j}^{2k - s}\phi^*\, \ipd{\sigma_j}^s \phi
      )
    )
    \\
    = &
    \int\limits_{\Omega} \dd[D]{x}\,
    \sqrt{- \det g}\,
    \qty(%
      \ap\,
      g^{\mu\nu}\, g^{\rho\sigma}\,
      \tensor{Ric}{_{\mu}_{\rho}}\,
      \qty(%
        c_{12}
        \phi^*\, \ipd{\nu\sigma}^2 \phi
        +
        c_{11}
        \ipd{\nu} \phi^*\, \ipd{\sigma} \phi
        +
        c_{10}
        \ipd{\nu\sigma}^2 \phi^*\, \phi
      )
    ),
  \end{split}
\end{equation}
where $\ap$ has been introduced after dimensional analysis and in order to have all adimensional $c$ factors.
Since only $\tensor{Ric}{_u_u}$ is non vanishing and it depends only on $u$,
the regularised d'Alembertian eigenmode problem now reads:
\begin{equation}
  \begin{split}
    -
    2 \ipd{u} \ipd{v} \phi_r
    & -
    \frac{u}{u^2 + \epsilon^2} \ipd{v} \phi_r
    +
    \frac{1}{\Delta^2 (u^2+ \epsilon^2)} \ipd{z}^2 \phi_r
    \\
    & +
    \finitesum{k}{1}{+\infty} \qty(\ap)^{2k-1}\,
    C_k\,
    \tensor{Ric}{_u_u}^k\,
    \ipd{v}^{2k} \phi
    +
    \ipd{i}^2 \phi_r
    -
    r\, \phi_r
    =
    0,
  \end{split}
\end{equation}
with $C_k = \finitesum{s}{0}{2k} (-1)^s\, c_{k s}$.
We can perform the usual Fourier transform and the function $B(u)$ becomes
\begin{equation}
  \begin{split}
    B(u)
    & =
    -
    (\vb{k}^2 + r)
    -
    i\, k_+\, \frac{u}{u^2 + \epsilon^2}
    -
    \frac{l^2}{\Delta^2}\, \frac{1}{u^2+\epsilon^2}
    \\
    & +
    \finitesum{k}{1}{+\infty}
    \qty(\ap)^{2k-1}\,
    C_k\,
    \qty(\frac{\epsilon^2}{(u^2 + \epsilon^2)^2})^k
    (-i k_+)^{2k}.
  \end{split}
\end{equation}
When $u = 0$ we have:
\begin{equation}
  B(0)
  \sim
  - \frac{l^2}{\Delta^2}\, \frac{1}{\epsilon^2}
  +
  \finitesum{k}{1}{+\infty}
  \qty(\ap)^{2k-1}\,
  C_k\,
  \frac{(-i k_+)^{2k}}{\epsilon^{2 k}}.
\end{equation}
Though the correction seems to lead to a cure for the divergence, ff we consider $\ap$ and $\epsilon^2$ uncorrelated we lose predictability.
However if $\ap \sim \epsilon^2$ as natural in string theory we do not solve the problem since 
\begin{equation}
  B(0)
  \stackrel{\ap \sim \epsilon^2}{\sim}
  - \frac{l^2}{\Delta^2}\, \frac{1}{\epsilon^2}
  +
  \finitesum{k}{1}{+\infty}
  C_k\,
  (-i k_+)^{2k}
  \epsilon^{2k - 2}
\end{equation}
and the curvature terms are no longer singular.
    

%%% TODO %%%
\subsection{A Hope from Twisted State Background}

It is clear from the previous discussion that the true problem is
associated with the dipole string and its charge neutral states
since the charged ones can be cured rather trivially by a Wilson line.

On the other side we know that the usual time-like orbifolds are well
defined because of a presence of a $B_{\mu\nu}$ background and this
field is sourced by strings.
So we can think of switching on such a background in the open string.
For open strings $B$ is equivalent to $F$ so we can consider what
happens to an open string in an electromagnetic background.

The choice of such a background is limited first of all
by the request that it must be an exact string solution, i.e. that it
satisfies the e.o.m derived from the DBI.
If a closed string winds the compact direction $z$
is coupled to $B_{z u}$, $B_{z v}$  and $B_{z i}$ but if we choose
\begin{equation}
\frac{1}{2\pi\ap}  B =  f(u) d u \wedge d z
  \label{eq:F_bck}
  .
  \end{equation}
then
  \begin{equation}
  \det(g+ 2\pi \ap f ) = \det(g)
  ,
  \end{equation}
therefore it is a solution of open string e.o.m. for any
$f(u,v,z,x^i)$.
 Suppose that the action for a real neutral scalar $\phi$ is given by
  (as the 2 tachyons -- 2 photons amplitude suggests)
  \begin{align}
    S_{\mbox{scalar kin}}
    =&
       \int_\Omega d^D x\,
       \sqrt{- \det g}
       \frac{1}{2}
       \Bigl(
       -g^{\alpha\beta}
       \ipd{\alpha}\phi\,
       \ipd{\beta} \phi
       -M^2  \phi^2
       + c_1
       \qty(\ap)^2\, \ipd{\mu} \phi \ipd{\nu}\phi
       f^{\mu\kappa} f^{\nu}_{~\kappa}
      \Bigr)
      \nonumber\\
    =
    &
      \int d^{D-3} \vb{x}\,
      \int d u\, \int d v\, \int_0^{2\pi} d z\,
      | \Delta u|
      \frac{1}{2}
      \Biggl(
      2\ipd{u} \phi\,\ipd{v} \phi\,
       \nonumber\\
       &
      -
         \frac{1}{(\Delta u )^2}
         (\ipd{z}  \phi)^2
         -
         (\vec \partial \phi)^2
         -
         M^2 \phi^2
         +c_1 \qty(\ap)^2
         \frac{1}{(\Delta u)^2} (\ipd{v} \phi)^2 f^2(u)
      \Biggr)
      ,
\end{align}
Performing the same steps as before we get
\begin{equation}
      B(u)
      =
      (-\vb{k}^2-r)
      +
      (-i k_+) \frac{1}{u}
      +
      \frac{(-l^2 + c_1 \qty(\ap)^2 f(u)^2 k_+^2)}{\Delta^2\, u^2}
      ,
\end{equation}
so even for a constant $f(u)=f_0$ we get a solution which solves the issues.
Notice however that the ``trivial'' solution $f=f_0 d u \wedge d z$ is
not so trivial in Minkowski coordinates $f=\frac{f_0}{x^-} d x^-
\wedge d x^2$.
Though appealing, the study of the string in the presence of this non trivial background needs a deeper analysis which goes beyond the scope of this paper.  

\subsection{Summary and Conclusions}

In the previous analysis it seems that string theory cannot do better than field theory when the latter does not exist, at least at the perturbative level where one deals with particles.
Moreover when spacetime becomes singular, the string massive modes are not anymore spectators.
Everything seems to suggest that issues with spacetime singularities are hidden into contact terms and interactions with massive states.
This would explain in an intuitive way why the eikonal approach to gravitational scattering works well: the eikonal is indeed concerned with three point massless interactions.
In fact it appears that the classical and quantum scattering on an electromagnetic wave~\cite{Jackiw:1992:ElectromagneticFieldsMassless} or gravitational wave \cite{tHooft:1987:GravitonDominanceUltrahighenergy} in \bo and \nbo are well behaved.
From this point of view the ACV approach~\cite{Soldate:1987:PartialwaveUnitarityClosedstring,Amati:1987:SuperstringCollisionsPlanckian} may be more sensible, especially when considering massive external states~\cite{Black:2012:HighEnergyString}.
Finally it seems that all issues are related with the Laplacian associated with the space-like subspace with vanishing volume at the singularity.
If there is a discrete zero eigenvalue the theory develops divergences. 

% vim: ft=tex
